% !TeX root = ../../main.tex
% =====================================================================
%  Modelado del neumático
%  Responsable:
%  Última revisión:
% =====================================================================
\chapter{Modelado del neumático}
\label{ch:modelado_neumatico}

% Las curvas de este capítulo usan la función mf(x,B,C,D,E) (Magic Formula),
% declarada una sola vez en Estilo/preambulo.tex.

\begin{objetivos}
  \item Saber qué es el Tire Test Consortium, qué se compra al entrar y qué
        contiene un fichero de ensayo.
  \item Convertir datos brutos de banco en curvas utilizables: agrupar por
        condición, quitar la histéresis del barrido y corregir los residuos.
  \item Escribir la Magic Formula, explicar qué controla cada coeficiente y
        deducir de ellos la rigidez de deriva.
  \item Ajustar el modelo con un procedimiento repetible y, sobre todo, saber
        detectar un ajuste que sale bien numéricamente y mal físicamente.
  \item Calibrar el factor de escala de agarre con el skidpad y explicar por qué
        el modelo sirve para comparar opciones y no para predecir tiempos.
\end{objetivos}

\fichasesion{90 minutos}{Capítulo \ref{ch:neumatico} y manejo básico de
Python o MATLAB}{Ordenador con \texttt{Codigo/dinamica/ajuste\_mf.py} y, si los
hay, datos de ensayo}

El capítulo anterior explica cómo se comporta un neumático. Este trata de
convertir ese comportamiento en \emph{números} que se puedan meter en una hoja de
cálculo o en un simulador. Sin ese paso, la dinámica vehicular se queda en
conversación de bar: se puede discutir si conviene más barra delante o detrás,
pero no se puede calcular.

\begin{clave}
El modelo de neumático es el \textbf{cimiento numérico de todo el bloque}. El
reparto de rigideces del capítulo \ref{ch:regimen_estacionario}, los diagramas
del \ref{ch:diagramas_ymd} y la simulación de vuelta del
\ref{ch:simulacion_vuelta} no son más fiables que él. Por eso conviene decidir
pronto cuánto esfuerzo merece: un modelo malo con el que se sabe convivir es
mucho mejor que uno bueno en el que se cree a ciegas.
\end{clave}

% ---------------------------------------------------------------------
\section{El Tire Test Consortium}
\label{sec:modelado_neumatico:01}
\index{Tire Test Consortium}\index{TTC}

\subsection{Qué es y qué se compra}

El \textbf{Formula SAE Tire Test Consortium} (TTC) es un consorcio de equipos
que financia entre todos campañas de ensayo de los neumáticos que se usan en la
competición \cite{ttc_data,kasprzak_ttc_2006}. Los ensayos se hacen en la
instalación de Calspan, sobre una banda continua recubierta de un abrasivo
(3M Safety Walk), con el neumático montado en un cabezal que controla carga
vertical, ángulo de deriva, caída, presión y velocidad.

Al entrar en el consorcio ---cuota por ronda y acuerdo de confidencialidad--- se
obtiene acceso a los datos de esa ronda. Los datos \textbf{no se pueden publicar
ni compartir}, lo cual tiene una consecuencia práctica para este libro y para el
informe de diseño: se pueden enseñar las conclusiones y el método, no los
ficheros.

\subsection{Tipos de ensayo}

\begin{center}
\begin{tabularx}{\textwidth}{@{}lXl@{}}
\toprule
\textbf{Ensayo} & \textbf{Qué se barre} & \textbf{Qué sale} \\
\midrule
\emph{Cornering} & Ángulo de deriva, a carga y caída fijas & $\Fy(\slipa)$, $\Mz(\slipa)$ \\
\emph{Drive/brake} & Deslizamiento longitudinal, sin deriva & $\Fx(\slipr)$ \\
Combinado & Deslizamiento a varios ángulos de deriva & Elipse de fricción \\
Relajación & Escalones de deriva & Longitud de relajación \\
\bottomrule
\end{tabularx}
\end{center}

La estructura de un ensayo de \emph{cornering} es la de la
\cref{fig:modelado:ensayo}: se fija una combinación de presión y caída, se
escalona la carga vertical y, dentro de cada escalón, se barre el ángulo de
deriva de un extremo a otro varias veces. El fichero es la concatenación de todos
esos tramos, y el primer trabajo consiste en volver a separarlos.

\begin{figure}[htbp]
\centering
\begin{tikzpicture}
\begin{axis}[cursopanel,
  xlabel={tiempo (\si{\second})}, ylabel={$\slipa$ (\si{\degree})},
  xmin=0, xmax=100, ymin=-14, ymax=14,
  xtick={0,25,...,100}, ytick={-12,-6,...,12},
  title={(a) Barrido de ángulo de deriva}]
\addplot[curvaA,domain=0:99.9,samples=500]{12*sin(360*x/12.5)};
\end{axis}
\end{tikzpicture}%
\hfill
\begin{tikzpicture}
\begin{axis}[cursopanel,
  xlabel={tiempo (\si{\second})}, ylabel={$\Fz$ (\si{\newton})},
  xmin=0, xmax=100, ymin=0, ymax=2300,
  xtick={0,25,...,100}, ytick={0,500,...,2000},
  title={(b) Escalones de carga vertical}]
\addplot[curvaB,domain=0:99.9,samples=500]{250+550*floor(x/25)};
\node[cursonota,anchor=north west] at (axis cs:4,2250)
     {cada escalón,\\un bloque de datos};
\end{axis}
\end{tikzpicture}
\caption[Estructura de un ensayo de neumático]{Estructura de un ensayo de
\emph{cornering}. El ángulo de deriva se barre en ciclos continuos (a) mientras
la carga vertical se mantiene por escalones (b); presión y caída son constantes
dentro de cada corrida. Reconstruir esta estructura a partir del fichero es el
primer paso del tratamiento de datos.}
\label{fig:modelado:ensayo}
\end{figure}

\subsection{Y si tu neumático no está}

Es el caso más frecuente en un equipo que empieza, y el nuestro. Las salidas
posibles, de mejor a peor:

\begin{enumerate}
  \item \textbf{Cambiar de neumático} a uno que sí esté ensayado. Es la opción
        que más se subestima: elegir el neumático \emph{por sus datos} es una
        decisión de ingeniería perfectamente defendible ante un juez.
  \item \textbf{Usar los datos de un neumático parecido} ---misma medida, mismo
        tipo de construcción, compuesto comparable--- como forma de curva, y
        ajustar la escala con medidas propias en pista (apartado
        \ref{sec:modelado_neumatico:05}).
  \item \textbf{Modelo mínimo con dos números}: un $\muy$ de pico conservador y
        una rigidez de deriva estimada. Da para calcular transferencia de carga
        y balance; no da para una simulación de vuelta creíble.
  \item \textbf{Ensayar uno mismo.} Un remolque instrumentado o un banco casero
        es un proyecto de temporada completo, pero es la única vía si se quiere
        independencia real.
\end{enumerate}

\begin{ojo}
Lo que \textbf{no} vale es coger la curva de un neumático cualquiera de internet,
meterla en el simulador y no volver a mencionarlo. El error no está en usar datos
ajenos ---todos empezamos así--- sino en no dejar escrito en el informe de qué
neumático son, en qué se parece al nuestro y qué incertidumbre añade eso a todo
lo que sale después. Un juez perdona un modelo pobre; no perdona que no sepas
cómo de pobre es.
\end{ojo}

% ---------------------------------------------------------------------
\section{Tratamiento de los datos brutos}
\label{sec:modelado_neumatico:02}
\index{tratamiento de datos}

\subsection{Los canales y sus trampas}

Un fichero de ensayo trae del orden de veinte canales. Los que se usan casi
siempre:

\begin{center}
\begin{tabularx}{\textwidth}{@{}llX@{}}
\toprule
\textbf{Canal} & \textbf{Magnitud} & \textbf{Para qué} \\
\midrule
\texttt{ET} & Tiempo transcurrido & Separar corridas y tramos \\
\texttt{SA} & Ángulo de deriva & Variable independiente lateral \\
\texttt{SR} & Deslizamiento longitudinal & Variable independiente longitudinal \\
\texttt{IA} & Ángulo de caída & Condición de ensayo \\
\texttt{FZ} & Carga vertical & Condición de ensayo; nunca es constante del todo \\
\texttt{FY}, \texttt{FX} & Fuerzas & Lo que se quiere modelar \\
\texttt{MZ} & Par autoalineante & Modelo de par y avance neumático \\
\texttt{P} & Presión & Condición de ensayo; sube al calentarse \\
\texttt{RL}, \texttt{RE} & Radio cargado y efectivo & Radio de rodadura, cinemática \\
\texttt{TSTI/TSTC/TSTO} & Temperaturas de superficie & Diagnóstico y descarte \\
\bottomrule
\end{tabularx}
\end{center}

\begin{ojo}
Antes de tocar nada, comprobar dos cosas y dejarlas escritas:

\begin{enumerate}
  \item \textbf{Unidades.} Hay ficheros en unidades imperiales (libras,
        pulgadas) y en SI. Un factor \num{4.448} olvidado da un coche que agarra
        el doble.
  \item \textbf{Convenio de signos.} Comprobar que $\Fy$ y $\slipa$ crecen a la
        vez. Si tienen signos opuestos, el fichero está en otro convenio: hay que
        cambiar el signo, no seguir adelante.
\end{enumerate}
\end{ojo}

\subsection{Los seis pasos}

\begin{enumerate}
  \item \textbf{Separar por condición}: presión, caída y carga nominal. Cada
        combinación es un conjunto de datos independiente.
  \item \textbf{Descartar el calentamiento}: el primer barrido de cada corrida se
        hace con el neumático todavía frío y no representa nada. Se tira.
  \item \textbf{Promediar barridos de ida y de vuelta.} Es lo que quita el bucle
        de histéresis de la \cref{fig:modelado:brutos}.
  \item \textbf{Normalizar por carga}: la carga real oscila alrededor de la
        nominal. O se trabaja con $\Fy/\Fz$, o se agrupa en ventanas estrechas.
  \item \textbf{Corregir los residuos}: a $\slipa = 0$ la fuerza lateral medida
        no es cero. Esa diferencia es la conicidad y el \emph{ply steer} del
        neumático, y en la Magic Formula se recoge con $\MFSh$ y $\MFSv$.
  \item \textbf{Agrupar y promediar} en intervalos de ángulo. El resultado es la
        curva limpia sobre la que se ajusta.
\end{enumerate}

\begin{figure}[htbp]
\centering
\begin{tikzpicture}
\begin{axis}[cursoplot,height=6.4cm,
  xlabel={$\slipa$ (\si{\degree})}, ylabel={$\Fy$ (\si{\newton})},
  xmin=-13, xmax=13, ymin=-2100, ymax=2100,
  xtick={-12,-8,...,12}, ytick={-2000,-1000,...,2000},
  legend pos=south east]
\addplot[fusalGris,line width=0.5pt,domain=0:12.5,samples=500,variable=\t]
  ( {12*sin(360*\t/12.5)},
    {mf( 12*sin(360*\t/12.5) - 0.9*cos(360*\t/12.5),
         0.3226,1.5,1620*(1-0.05*\t/12.5),0.3 )} );
  \addlegendentry{dato bruto de un barrido}
\addplot[curvaB,domain=-13:13,samples=200]{mf(x,0.3226,1.5,1580,0.3)};
  \addlegendentry{curva promediada}
\node[cursonota,anchor=north west] at (axis cs:-12.5,2000)
     {el bucle es transitorio térmico y de\\relajación, no comportamiento};
\end{axis}
\end{tikzpicture}
\caption[Dato bruto frente a curva promediada]{Un barrido completo de ángulo de
deriva tal como sale del banco (línea fina) y la curva que se obtiene al
promediar ida y vuelta (línea gruesa). El bucle no es una propiedad del
neumático: es el retraso con que responde a un ángulo que está cambiando, más el
calentamiento a lo largo del barrido. Ajustar el modelo sobre el dato bruto sin
promediar mete ese bucle dentro de los coeficientes. Curvas sintéticas.}
\label{fig:modelado:brutos}
\end{figure}

% ---------------------------------------------------------------------
\section{La Magic Formula de Pacejka}
\label{sec:modelado_neumatico:03}
\index{Magic Formula}\index{Pacejka}

\subsection{La fórmula}

La Magic Formula \cite{bakker_tiremodel_1987,pacejka_tirevehicle_2012} es una
expresión empírica ---no se deduce de la física del contacto, se ajusta a los
datos--- que reproduce muy bien la forma de las curvas del neumático con muy
pocos parámetros:

\begin{equation}
  y(x) = \MFD \sin\Bigl(\MFC \arctan\bigl(\MFB x -
         \MFE(\MFB x - \arctan \MFB x)\bigr)\Bigr) + \MFSv,
  \qquad x = X + \MFSh
  \label{eq:magic_formula}
\end{equation}

donde $X$ es $\slipa$ o $\slipr$ e $y$ es $\Fy$, $\Fx$ o $\Mz$. El nombre viene
de que nadie sabe explicar por qué funciona tan bien; lo que sí se puede explicar
es qué hace cada coeficiente (\cref{fig:modelado:coeficientes}):

\begin{center}
\begin{tabularx}{\textwidth}{@{}llX@{}}
\toprule
\textbf{Coef.} & \textbf{Nombre} & \textbf{Qué controla} \\
\midrule
$\MFD$ & Valor de pico & La altura de la curva: $\MFD = \muy \Fz$ \\
$\MFC$ & Factor de forma & Cuánto cae la curva después del pico \\
$\MFB$ & Factor de rigidez & Dónde está el pico, y con él la pendiente inicial \\
$\MFE$ & Factor de curvatura & La forma justo alrededor del pico \\
$\MFSh$, $\MFSv$ & Desplazamientos & Conicidad, \emph{ply steer}, resistencia a la rodadura \\
\bottomrule
\end{tabularx}
\end{center}

\begin{clave}
La relación más útil de todo el capítulo:

\begin{equation}
  \Calpha = \MFB\,\MFC\,\MFD
  \label{eq:bcd}
\end{equation}

El producto de los tres primeros coeficientes es la \textbf{rigidez de deriva},
es decir, la pendiente en el origen. Sirve para dos cosas: comprobar de un
vistazo si un ajuste es razonable ---si $\MFB\MFC\MFD$ da
\SI{4000}{\newton\per\degree}, el ajuste está mal--- y alimentar directamente los
modelos lineales de balance del capítulo \ref{ch:regimen_estacionario}, que no
necesitan la curva entera.
\end{clave}

\begin{figure}[htbp]
\centering
\begin{tikzpicture}
\begin{axis}[cursopanel,height=5.2cm,
  xlabel={$\slipa$ (\si{\degree})}, ylabel={$\Fy$ (\si{\newton})},
  xmin=0, xmax=12, ymin=0, ymax=2400, ytick={0,500,...,2000},
  title={$\MFD$: altura del pico}]
\addplot[curvaA,domain=0:12,samples=100]{mf(x,0.3226,1.5,1200,0.3)};
\addplot[curvaB,domain=0:12,samples=100]{mf(x,0.3226,1.5,1620,0.3)};
\addplot[curvaC,domain=0:12,samples=100]{mf(x,0.3226,1.5,2040,0.3)};
\end{axis}
\end{tikzpicture}%
\hfill
\begin{tikzpicture}
\begin{axis}[cursopanel,height=5.2cm,
  xlabel={$\slipa$ (\si{\degree})}, ylabel={$\Fy$ (\si{\newton})},
  xmin=0, xmax=12, ymin=0, ymax=2400, ytick={0,500,...,2000},
  title={$\MFB$: posición del pico}]
\addplot[curvaA,domain=0:12,samples=100]{mf(x,0.2200,1.5,1620,0.3)};
\addplot[curvaB,domain=0:12,samples=100]{mf(x,0.3226,1.5,1620,0.3)};
\addplot[curvaC,domain=0:12,samples=100]{mf(x,0.5000,1.5,1620,0.3)};
\end{axis}
\end{tikzpicture}

\vspace{3mm}
\begin{tikzpicture}
\begin{axis}[cursopanel,height=5.2cm,
  xlabel={$\slipa$ (\si{\degree})}, ylabel={$\Fy$ (\si{\newton})},
  xmin=0, xmax=12, ymin=0, ymax=2400, ytick={0,500,...,2000},
  title={$\MFC$: caída tras el pico}]
\addplot[curvaA,domain=0:12,samples=100]{mf(x,0.3226,1.25,1620,0.3)};
\addplot[curvaB,domain=0:12,samples=100]{mf(x,0.3226,1.5,1620,0.3)};
\addplot[curvaC,domain=0:12,samples=100]{mf(x,0.3226,1.8,1620,0.3)};
\end{axis}
\end{tikzpicture}%
\hfill
\begin{tikzpicture}
\begin{axis}[cursopanel,height=5.2cm,
  xlabel={$\slipa$ (\si{\degree})}, ylabel={$\Fy$ (\si{\newton})},
  xmin=0, xmax=12, ymin=0, ymax=2400, ytick={0,500,...,2000},
  title={$\MFE$: forma en el pico}]
\addplot[curvaA,domain=0:12,samples=100]{mf(x,0.3226,1.5,1620,-1.0)};
\addplot[curvaB,domain=0:12,samples=100]{mf(x,0.3226,1.5,1620,0.3)};
\addplot[curvaC,domain=0:12,samples=100]{mf(x,0.3226,1.5,1620,0.9)};
\end{axis}
\end{tikzpicture}
\caption[Efecto de los coeficientes de la Magic Formula]{Qué hace cada
coeficiente de la Magic Formula. En cada panel se varía uno y se dejan fijos los
otros tres; la curva roja es la misma en los cuatro. Entender esta figura es lo
que permite poner semillas razonables al ajuste y, sobre todo, darse cuenta de
cuándo un ajuste ha convergido a un sitio absurdo.}
\label{fig:modelado:coeficientes}
\end{figure}

\subsection{Dependencia con la carga y con la caída}

Un juego de cuatro coeficientes solo vale para una carga vertical. Como el coche
trabaja a cargas que cambian continuamente, la Magic Formula completa hace que
los coeficientes dependan de la variación normalizada de carga

\begin{equation}
  \dfz = \frac{\Fz - \Fzo}{\Fzo}
  \label{eq:dfz}
\end{equation}

mediante polinomios. Por ejemplo, el valor de pico y la rigidez de deriva se
escriben habitualmente como

\begin{equation}
  \MFD = (p_{D1} + p_{D2}\,\dfz)\,\Fz, \qquad
  \Calpha = p_{K1}\,\Fzo \sin\!\left(2\arctan\frac{\Fz}{p_{K2}\,\Fzo}\right)
  \label{eq:dependencia_carga}
\end{equation}

y la caída entra como términos adicionales en $\MFSh$, $\MFSv$ y $\MFD$. Ahí es
donde la fórmula deja de ser elegante: el modelo completo tiene más de cien
coeficientes. Para Formula Student casi nunca hacen falta todos.

\begin{center}
\begin{tabularx}{\textwidth}{@{}lXl@{}}
\toprule
\textbf{Nivel de modelo} & \textbf{Cuándo basta} & \textbf{Parámetros} \\
\midrule
Lineal, $\Fy = \Calpha \slipa$ & Balance y gradiente de subviraje & 1 por eje \\
MF por carga, interpolando & Transferencia de carga, diagramas & 4--6 por carga \\
MF con dependencia de $\dfz$ & Simulación de vuelta & 10--20 \\
MF completa (fichero \texttt{.tir}) & Simulación multicuerpo & \num{100}+ \\
\bottomrule
\end{tabularx}
\end{center}

\begin{ojo}
La tentación es ir directamente al modelo más completo. Es un error de orden: el
modelo se elige por la \emph{pregunta} que hay que contestar. Un modelo lineal
contesta bien «¿de qué lado va el balance?» en diez minutos, y esa suele ser la
pregunta que de verdad hay encima de la mesa. Ajustar cien coeficientes a datos
de un neumático que ni siquiera es el nuestro es trabajo que parece ingeniería y
no lo es.
\end{ojo}

% ---------------------------------------------------------------------
\section{Ajuste y validación del modelo}
\label{sec:modelado_neumatico:04}
\index{ajuste de curvas}\index{validación}

\subsection{Procedimiento}

El ajuste es un problema de mínimos cuadrados no lineales. Lo que lo hace fácil o
imposible no es el algoritmo, sino cómo se le plantea:

\begin{enumerate}
  \item \textbf{Ajustar carga por carga, no todo a la vez.} Primero un juego de
        coeficientes por cada escalón de carga; después, la dependencia con la
        carga sobre esos resultados. Un ajuste global de veinte parámetros a la
        vez casi nunca converge a algo con sentido.
  \item \textbf{Dar semillas razonables}: $\MFD \approx$ el máximo de los datos,
        $\MFC \approx \num{1.5}$ en lateral y $\approx \num{1.65}$ en
        longitudinal, $\MFB$ tal que el pico caiga donde cae en los datos,
        $\MFE \approx 0$.
  \item \textbf{Poner límites} a cada parámetro. Sin límites, el optimizador
        encuentra soluciones con $\MFD$ negativo que ajustan estupendamente y no
        significan nada.
  \item \textbf{Ponderar la zona lineal.} Hay muchos más puntos cerca del pico
        que cerca del origen, y el ajuste sacrifica la pendiente inicial ---que
        es justo lo que gobierna el comportamiento normal del coche--- para ganar
        un poco en el pico.
\end{enumerate}

El programa \texttt{ajuste\_mf.py} implementa este procedimiento. Genera datos
sintéticos si no se le pasa fichero, de modo que el flujo completo se puede
probar antes de tener datos reales:

\lstinputlisting[language=Python, firstline=43, lastline=60,
  caption={Modelo y rigidez de deriva (\texttt{Codigo/dinamica/ajuste\_mf.py})},
  label={lst:mf_modelo}]{dinamica/ajuste_mf.py}

\lstinputlisting[language=Python, firstline=111, lastline=145,
  caption={Ajuste de un escalón de carga y comprobaciones de sensatez},
  label={lst:mf_ajuste}]{dinamica/ajuste_mf.py}

\subsection{Cómo saber si el ajuste es bueno}

Un coeficiente de determinación alto no dice gran cosa: la Magic Formula tiene
tanta libertad que ajusta casi cualquier nube de puntos. Lo que hay que mirar es
otra cosa (\cref{fig:modelado:ajuste}):

\begin{itemize}
  \item \textbf{Los residuos, no el error medio.} Si los residuos tienen
        estructura ---todos positivos en una zona y negativos en otra--- el
        modelo no está capturando la forma, por bueno que sea el $R^2$.
  \item \textbf{$\MFB\MFC\MFD$ frente a la pendiente medida} en los primeros
        grados. Si no coinciden, el ajuste ha sacrificado la zona lineal.
  \item \textbf{$\MFD/\Fz$ dentro de lo posible.} Un $\muy$ de \num{3.5} no es un
        neumático: es un error de unidades.
  \item \textbf{Una condición reservada.} Ajustar sin usar una de las cargas y
        comprobar después qué predice el modelo para ella. Es la única
        comprobación que de verdad mide capacidad predictiva.
\end{itemize}

\begin{figure}[htbp]
\centering
\begin{tikzpicture}
\begin{axis}[cursoplot,height=5.4cm,
  ylabel={$\Fy$ (\si{\newton})},
  xmin=0, xmax=12.5, ymin=0, ymax=1900,
  xtick={0,2,...,12}, xticklabels={},
  ytick={0,500,...,1500},
  legend pos=south east]
\addplot[curvaB,domain=0:12.5,samples=150]{mf(x,0.3226,1.5,1620,0.3)};
  \addlegendentry{Magic Formula ajustada}
\addplot[only marks,mark=*,mark size=1.3pt,fusalPrim] coordinates {
  (0,6) (1,740) (2,1187) (3,1459) (4,1548) (5,1618) (6,1607)
  (7,1631) (8,1589) (9,1598) (10,1581) (11,1547) (12,1544)};
  \addlegendentry{datos promediados}
\end{axis}
\end{tikzpicture}

\begin{tikzpicture}
\begin{axis}[cursoplot,height=2.8cm,
  xlabel={$\slipa$ (\si{\degree})}, ylabel={residuo (\si{\newton})},
  xmin=0, xmax=12.5, ymin=-30, ymax=30,
  xtick={0,2,...,12}, ytick={-20,0,20}]
\addplot[only marks,mark=*,mark size=1.3pt,fusalSec] coordinates {
  (0,6) (1,15) (2,-15) (3,11) (4,-13) (5,11) (6,-13)
  (7,15) (8,-16) (9,8) (10,8) (11,-9) (12,5)};
\end{axis}
\end{tikzpicture}
\caption[Ajuste de la Magic Formula y residuos]{Ajuste de la Magic Formula sobre
datos promediados, con el panel de residuos debajo. Lo que se busca en el panel
inferior es \emph{ruido sin estructura}: si los residuos formaran una curva
---por ejemplo, todos negativos entre \num{2} y \SI{5}{\degree}--- el modelo
estaría fallando en la forma, y habría que revisar $\MFE$ o el planteamiento del
ajuste en vez de aceptar el resultado porque el error medio sea pequeño. Datos
sintéticos.}
\label{fig:modelado:ajuste}
\end{figure}

\begin{ojo}
\textbf{Nunca extrapolar.} Fuera del rango de ángulos, cargas, presiones y caídas
ensayado, la Magic Formula no es un modelo: es un polinomio suelto haciendo lo
que le apetece. El simulador acabará pidiendo cargas de \SI{3000}{\newton} en un
bordillo aunque el ensayo llegara a \num{1600}, y devolverá un número sin avisar.
Hay que \textbf{saturar el modelo} en los extremos del rango ensayado y contar
cuántas veces se ha llegado a esa saturación: si es a menudo, el resultado de la
simulación no vale.
\end{ojo}

% ---------------------------------------------------------------------
\section{Factores de escala y agarre real en pista}
\label{sec:modelado_neumatico:05}
\index{factor de escala}\index{skidpad}

\subsection{El problema de la banda de ensayo}

Los datos del TTC se toman sobre 3M Safety Walk, un abrasivo mucho más agresivo
que el asfalto. Un neumático que en el banco da $\muy = \num{2.6}$ en pista da
entre \num{1.4} y \num{1.7}. Si esos datos se meten sin corregir en el simulador,
sale un coche que da \SI{2.5}{\gacc} laterales y que se diseña ---en frenos, en
aerodinámica, en rigidez de suspensión--- para unas cargas que nunca va a ver.

La corrección habitual es un factor de escala global sobre el pico:

\begin{equation}
  \MFD^{\text{pista}} = \lamu \, \MFD^{\text{banco}}
  \label{eq:factor_escala}
\end{equation}

con $\lamu$ habitualmente entre \num{0.55} y \num{0.75}. Es una corrección
grosera ---escala el pico pero no la rigidez de deriva, que en
realidad también cambia--- y es, con diferencia, la mayor fuente de incertidumbre
de todo el bloque de dinámica. Por eso no puede quedarse como un número heredado
en un script: hay que medirlo.

\subsection{Calibrar con el skidpad}

El skidpad es un ensayo de agarre lateral casi puro y sale gratis: se corre en
cada competición y se puede repetir en cualquier aparcamiento con conos. Para un
radio de trayectoria $R$ y un tiempo de vuelta $t$:

\begin{equation}
  \ay = \frac{v^2}{R} = \frac{(2\pi R/t)^2}{R}
  \quad\Longrightarrow\quad
  \muy^{\text{efectivo}} = \frac{4\pi^2 R}{g\,t^2}
  \label{eq:skidpad}
\end{equation}

Con el radio nominal de \SI{9.125}{\meter}, un tiempo de \SI{5.2}{\second} da un
$\muy$ efectivo de \num{1.36} (\cref{fig:modelado:escala}b). Ese número se
compara con el que da el modelo, y la diferencia \emph{es} el factor de escala.

\begin{clave}
El $\muy$ que sale del skidpad es un \textbf{coeficiente efectivo del coche
entero}, no el pico del neumático: incluye la pérdida por transferencia de carga,
el hecho de que las cuatro ruedas no trabajan en su óptimo a la vez y que el
piloto no está exactamente en el límite. Por eso es siempre \emph{menor} que el
pico del neumático, y por eso es exactamente el número que hay que reproducir con
el modelo completo del coche: si tu simulación del skidpad da el tiempo medido,
el modelo de neumático y el de transferencia de carga son coherentes entre sí.

A la velocidad del skidpad ---unos \SI{40}{\kilo\meter\per\hour}--- la carga
aerodinámica es despreciable, lo que convierte este ensayo en la medida de agarre
más limpia que un equipo puede hacer sin instrumentación.
\end{clave}

\begin{figure}[htbp]
\centering
\begin{tikzpicture}
\begin{axis}[cursopanel,
  xlabel={$\slipa$ (\si{\degree})}, ylabel={$\Fy$ (\si{\newton})},
  xmin=0, xmax=12, ymin=0, ymax=2900,
  xtick={0,2,...,12}, ytick={0,500,...,2500},
  title={(a) Escalado del pico}, legend pos=south east]
\addplot[curvaA,domain=0:12,samples=120]{mf(x,0.3226,1.5,2600,0.3)};
  \addlegendentry{banco, $\muy=\num{2.6}$}
\addplot[curvaB,domain=0:12,samples=120]{mf(x,0.3226,1.5,1560,0.3)};
  \addlegendentry{pista, $\lamu=\num{0.6}$}
\end{axis}
\end{tikzpicture}%
\hfill
\begin{tikzpicture}
\begin{axis}[cursopanel,
  xlabel={$\muy$ efectivo}, ylabel={tiempo de skidpad (\si{\second})},
  xmin=0.9, xmax=2.0, ymin=4.2, ymax=6.6,
  xtick={1.0,1.2,...,2.0}, ytick={4.5,5.0,...,6.5},
  title={(b) Calibración con el skidpad}]
\addplot[curvaC,domain=0.9:2.0,samples=120]{6.0605/sqrt(x)};
\draw[cursomarca] (axis cs:0.9,5.2) -- (axis cs:1.358,5.2)
                  -- (axis cs:1.358,4.2);
\node[cursonota,anchor=south west] at (axis cs:1.40,5.3)
     {\SI{5.2}{\second} $\Rightarrow$\\$\muy \approx \num{1.36}$};
\end{axis}
\end{tikzpicture}
\caption[Escalado del modelo y calibración con el skidpad]{(a) La curva de banco
y la misma curva escalada al agarre de pista: el factor $\lamu$ recorta la fuerza
disponible casi a la mitad. (b) Tiempo de vuelta de skidpad frente al coeficiente
de agarre efectivo, según \eqref{eq:skidpad} con $R=\SI{9.125}{\meter}$. La curva
se lee al revés de como se dibuja: se mide el tiempo y se deduce el agarre real
del coche.}
\label{fig:modelado:escala}
\end{figure}

\subsection{Qué aguanta y qué no el modelo escalado}

Merece la pena tener claro qué conclusiones sobreviven a un error del
\SI{10}{\percent} en $\lamu$ y cuáles no:

\begin{center}
\begin{tabularx}{\textwidth}{@{}Xl@{}}
\toprule
\textbf{Pregunta} & \textbf{¿Fiable?} \\
\midrule
¿Qué reparto de barras equilibra el coche? & Sí, muy poco sensible \\
¿Conviene más rigidez delante o detrás? & Sí \\
¿Cuánto gano quitando \SI{10}{\kilogram}? & Sí, en términos relativos \\
¿Qué relación de transmisión pongo? & Sí \\
¿Qué tiempo voy a hacer en el autocross? & No \\
¿Cuántas $g$ laterales da el coche? & No: escala directamente con $\lamu$ \\
¿Qué cargas uso para dimensionar la suspensión? & No sin un margen explícito \\
\bottomrule
\end{tabularx}
\end{center}

\begin{clave}
Un modelo de neumático sirve para \textbf{ordenar opciones}, no para predecir
resultados. Las respuestas relativas ---esto es mejor que aquello, y por este
margen--- se mantienen aunque el factor de escala esté equivocado, porque el
error afecta igual a las dos alternativas. Las respuestas absolutas ---este
tiempo, estas $g$, estas cargas--- se van con él.

Esa es también la forma correcta de contarlo en el design event: no «nuestro
coche da \SI{1.45}{\gacc}», sino «con el modelo escalado a partir del skidpad,
esta configuración da un \SI{4}{\percent} más de agarre de eje que la
alternativa, y así es como lo hemos comprobado».
\end{clave}

% ---------------------------------------------------------------------
%  Cierre del capítulo
% ---------------------------------------------------------------------

\begin{caso}[Modelar sin datos]
No estamos en el TTC y el Pirelli \texttt{185/40 R13} que montamos no está
ensayado en ninguna ronda, así que a día de hoy \textbf{no tenemos modelo de
neumático}. Lo que hay es un $\muy$ conservador metido a mano en las hojas de
cálculo de transferencia de carga y de casos de carga.

Eso condiciona todo el bloque, y conviene decirlo tal cual: los diagramas del
capítulo \ref{ch:diagramas_ymd} y la simulación de vuelta del
\ref{ch:simulacion_vuelta} no se pueden cerrar con rigor mientras esto siga así.
Nuestras prioridades, en orden:

\begin{enumerate}
  \item \textbf{Medir el skidpad} en cuanto el coche ruede. Es gratis, no
        necesita instrumentación y da directamente el $\muy$ efectivo con el que
        estamos rodando, que hoy es literalmente un número supuesto.
  \item \textbf{Decidir con datos el neumático de la temporada siguiente}: si el
        candidato está en el TTC, entrar en el consorcio pasa a ser una compra
        justificada y no un capricho.
  \item \textbf{Montar el flujo de ajuste antes de tener los datos}, con
        \texttt{ajuste\_mf.py} y sus datos sintéticos, para que el día que
        lleguen no se pierda un mes aprendiendo a leerlos.
\end{enumerate}

\pendiente{Cerrar este recuadro con: tiempo de skidpad medido, $\muy$ efectivo
deducido y el valor de $\muy$ que se estaba usando antes en las hojas de cálculo.
La diferencia entre los dos últimos números es lo que este capítulo tiene que
enseñar al equipo.}
\end{caso}

\begin{ojo}
\begin{itemize}
  \item Meter los datos del TTC sin escalar y obtener \SI{2.5}{\gacc} laterales
        sin que a nadie le extrañe.
  \item Ajustar los cien coeficientes del modelo completo con datos de cuatro
        cargas: sobra modelo y falta información.
  \item Aceptar un ajuste porque el $R^2$ es \num{0.99}, sin mirar los residuos
        ni comprobar $\MFB\MFC\MFD$.
  \item Extrapolar fuera del rango ensayado y no enterarse, porque el modelo no
        avisa.
  \item Perder el rastro de qué versión del modelo produjo qué resultado. Los
        coeficientes van en un fichero con fecha y con el mismo control de
        versiones que el código.
  \item Confundir el $\muy$ del neumático con el $\muy$ efectivo del coche: el
        segundo es siempre menor, y es el único que se mide.
  \item Repetir el ajuste a mano cada vez. Si no es un script que se ejecuta
        entero de principio a fin, dentro de dos temporadas nadie sabrá
        reproducirlo.
\end{itemize}
\end{ojo}

\begin{entregable}
\begin{enumerate}
  \item El \textbf{fichero de coeficientes} del neumático de la temporada, con
        fecha, origen de los datos, rango de validez (ángulos, cargas, presiones,
        caídas) y factor de escala aplicado. Es el fichero del que dependen todos
        los cálculos del bloque.
  \item El \textbf{script de ajuste} en el repositorio, ejecutable de principio a
        fin sobre los datos brutos, sin pasos manuales.
  \item La \textbf{página de validación} del informe de diseño: gráfica de ajuste
        con residuos, comparación de $\MFB\MFC\MFD$ con la pendiente medida y
        calibración de $\lamu$ con el skidpad.
  \item El \textbf{registro de skidpad}: tiempo, radio usado, presiones,
        temperaturas y $\muy$ efectivo deducido, una fila por sesión.
\end{enumerate}
\end{entregable}

\begin{juez}
\begin{enumerate}
  \item ¿De dónde salen los datos de vuestro neumático?
  \item ¿Qué factor de escala aplicáis y cómo lo habéis obtenido?
  \item ¿Qué $\muy$ efectivo medisteis en el skidpad y cuánto se parece al que
        predice vuestro modelo?
  \item Enseñadme el ajuste y los residuos.
  \item ¿Qué hace vuestro simulador cuando la carga se sale del rango ensayado?
  \item ¿Cuánto cambiarían vuestras conclusiones si el factor de escala estuviera
        equivocado un \SI{10}{\percent}?
\end{enumerate}
\end{juez}
